% --------------------------
% ---- DECLARE PACKAGES ----
% --------------------------
\documentclass[a4paper, oneside, openright]{book}

% PDF/A-2b
\usepackage[a-2b]{pdfx}

\usepackage[T1]{fontenc} % Font encoding, T1 = it
\usepackage{lmodern}
\usepackage[utf8]{inputenc} % Input encoding
\usepackage[english,italian]{babel} % Lingua principale italiana, con parti in inglese
\usepackage{graphicx} % Immagini esterne

\usepackage{setspace}
\onehalfspacing % interlinea 1.5

\usepackage[parfill]{parskip} % Disabilita l'indentazione dopo essere andati a capo
\usepackage{titlesec}
\usepackage{booktabs}
\usepackage{longtable}
\usepackage{tabularx}
\usepackage{enumitem}
\usepackage{amsmath}
\usepackage{amssymb}

% Colonne di tabella a bandiera
\newcolumntype{L}{>{\raggedright\arraybackslash}X}
\newcolumntype{P}[1]{>{\raggedright\arraybackslash}p{#1}}
\newcolumntype{B}[1]{>{\raggedright\arraybackslash\bfseries}p{#1}}

% Percorso immagini
\graphicspath{{../Immagini/}{../Immagini/Sessioni/}{../Immagini/Screenshot/}}

% Diagrammi TikZ
\usepackage{tikz}
\usepackage{pgfplots}
\pgfplotsset{compat=1.18}
\usetikzlibrary{shapes.geometric, shapes.multipart, arrows, arrows.meta,
                positioning, fit, backgrounds, calc, chains, automata,
                decorations.pathreplacing}
\tikzstyle{startstop} = [rectangle, rounded corners, minimum width=3cm, minimum height=1cm,text centered, draw=black, fill=red!30]
\tikzstyle{io} = [trapezium, trapezium left angle=70, trapezium right angle=110, minimum width=3cm, minimum height=1cm, text centered, draw=black, fill=blue!30]
\tikzstyle{process} = [rectangle, minimum width=3cm, minimum height=1cm, text centered, draw=black, fill=orange!30]
\tikzstyle{decision} = [diamond, minimum width=3cm, minimum height=1cm, text centered, draw=black, fill=green!30]
\tikzstyle{arrow} = [thick,->,>=stealth]

% Stili per i diagrammi
\tikzset{
  layer/.style   = {rectangle, rounded corners=3pt, draw=black!70, thick,
                    align=center, inner sep=6pt, minimum height=1cm},
  comp/.style    = {rectangle, rounded corners=2pt, draw=black!60,
                    fill=white, align=center, font=\footnotesize,
                    inner sep=4pt, minimum height=0.8cm},
  extern/.style  = {rectangle, dashed, draw=black!60, align=center,
                    font=\footnotesize, inner sep=4pt},
  ucase/.style   = {ellipse, draw=black!70, align=center,
                    font=\footnotesize, inner sep=2pt, minimum width=2.6cm},
  flow/.style    = {-{Latex[length=2mm]}, thick},
  flowdash/.style= {-{Latex[length=2mm]}, thick, dashed},
  % Stili per i diagrammi
  tile/.style    = {rectangle, draw=black!55, fill=black!2, thin},
  cellgrid/.style= {draw=black!35, thin, dash pattern=on 1.2pt off 1.2pt},
  pin/.style     = {circle, fill=black!75, draw=none, inner sep=0pt,
                    minimum size=3pt},
  bubble/.style  = {circle, draw=black!65, fill=black!12, thick,
                    inner sep=0pt, font=\scriptsize},
  tick/.style    = {black!70, thin},
  band/.style    = {rectangle, draw=black!55, fill=black!6, rounded corners=1pt}
}

% Captions
\usepackage[font=scriptsize, skip=5pt, justification=centering]{caption}

\usepackage[bottom]{footmisc}
\usepackage{hyperref}
\usepackage{subcaption}

% Sperimentali
\usepackage[immediate]{silence} % Necessario per \WarningFilter
\WarningFilter{microtype}{Unable to apply patch `footnote'}
\usepackage{microtype}

\usepackage[top=3.5cm,bottom=3.5cm,left=3.5cm,right=3.5cm]{geometry}

% ------------------------
% ---- DOCUMENT SETUP ----
% ------------------------
\pagestyle{plain}
\raggedbottom

% Titolo del capitolo unico, senza prefisso e senza numerazione
\titleformat{\chapter}[display]
{\normalfont\fontsize{20pt}{26pt}\selectfont\bfseries}{}{0pt}{}
\titlespacing*{\chapter}{0pt}{0pt}{18pt}

\hypersetup{
    colorlinks,
    citecolor=black,
    filecolor=black,
    linkcolor=black,
    urlcolor=black
}

\renewcommand{\footnotesize}{\fontsize{11pt}{13pt}\selectfont} % Imposta la dimensione del testo delle footnotes
\setlength{\footnotesep}{0.5cm} % Imposta lo spazio fra le singole footnotes
\setlength{\skip\footins}{1.5cm} % Imposta lo spazio fra il corpo e le footnotes

\DeclareUnicodeCharacter{02BC}{}

% ------------------------
% ---- DOCUMENT START ----
% ------------------------
\begin{document}

\chapter*{Riassunto tesi Gabriele Caretti 756564}
\label{ch:riassunto}

La dermatite atopica e la psoriasi sono patologie infiammatorie croniche che non guariscono.
Chi ne è affetto ha bisogno di un percorso di cure, con controlli periodici e trattamenti che solo strutture autorizzate
possono prescrivere.
L'utente quindi si pone una domanda, cioè dove si trovi il centro più vicino che segue la propria
patologia.
Esiste una soluzione, ma frammentata, sotto forma di elenchi e mappe web organizzati per regione o provincia,
che richiedono una connessione attiva.
Il lavoro oggetto della tesi ha realizzato un'applicazione mobile cross-platform, in Flutter, che rende interrogabili per
posizione dati già pubblici e raccoglie da due portali i centri dermatologici dedicati alle due patologie,
li normalizza e li mostra su una mappa interattiva con clustering dinamico dei marker, filtri per patologia e per
raggio, ricerca per indirizzo e centri personalizzati inseriti dall'utente, delegando le azioni successive
come il calcolo del percorso a programmi esterni installati sul dispositivo.
L'applicazione non formula diagnosi, non prevede autenticazione né backend dedicato e non trasmette all'esterno nessuna
informazione dell'utente.

%Le scelte più condizionanti riguardano i dati e la cartografia.
Non essendoci un proprio servizio per i dati, la responsabilità di rendere omogenee due fonti eterogenee è a carico del client.
Il modulo è organizzato in tre livelli, dei quali solamente il livello di accesso conosce la forma originaria dei record, e la
normalizzazione avviene per catene di alias, offrendo grande tolleranza ai cambiamenti che potrebbero avvenire a monte sulle fonti. %, più tolleranti ai cambiamenti di formato di una classe di dominio tipizzata.
Come provider delle mappe si è scelto OpenStreetMap con \texttt{flutter\_map} anziché Google Maps o Mapbox, che avrebbero
richiesto una chiave associata a un profilo di fatturazione, e i marker sono widget Flutter.
I dati vengono conservati su tre archivi distinti, di cui uno riservato ai centri inseriti a mano, %che nessun aggiornamento può sovrascrivere,
così che l'applicazione resti utilizzabile anche senza connessione.

A causa dell'utilizzo di OpenStreetMap rispetto a un provider commerciale, l'algoritmo di clustering è stato
sviluppato internamente.
Lo spazio delle coordinate è suddiviso in celle quadrate il cui lato dimezza a ogni incremento
dello zoom, seguendo lo schema \emph{slippy map}, così che la loro dimensione sullo schermo resti costante e i centri
posizionati nelle aree urbane restino leggibili.
%La verifica si è limitata a ciò che è ripetibile senza dispositivo e senza rete.
Sono stati creati dodici unit test sulla serializzazione, sulla persistenza locale dei dati e sul calcolo geodetico,
completati da verifiche manuali su Android.
%I limiti principali sono due: la cache degli elenchi non ha una scadenza e l'applicazione dipende da elenchi pubblici di
%terzi, sui quali non si ha controllo.
Viene inoltre proposta una possibile futura direzione di sviluppo, ossia l'introduzione di un livello intermedio di normalizzazione
fra le sorgenti e l'applicazione, che renderebbe superflue le catene di alias, permetterebbe di riconoscere i cambiamenti
di formato sul dispositivo e ridurrebbe le richieste ai servizi pubblici.

% ----------------------
% ---- DOCUMENT END ----
% ----------------------
\end{document}
